Situémonos ahora bajo las condiciones de~\Ref{X.2.4}, donde se tiene un
homomorfismo de especialización
$$
\pi_1(\overline X_1,\overline a_1)\to\pi_1(\overline X_0,\overline
a_0)
$$
que es \emph{sobreyectivo}, relativo a un morfismo propio y separable
$f\colon X\to Y$. Queremos precisar el núcleo de este homomorfismo. Procediendo
como en la demostración de~\Ref{X.2.4}, se ve que en esta cuestión siempre
podemos suponer que~$Y$ es el espectro de un \emph{anillo de valoración
discreta~$V$, completo y de cuerpo residual algebraicamente cerrado} (pues
siempre puede encontrarse un anillo semejante y un morfismo de su espectro~$Y'$
en~$Y$ cuya imagen sea~$\{y_0,y_1\}$). Se tiene entonces
$X_0=\overline X_0$,~$\kres(y_0)=k=$ cuerpo residual de~$V$ y
$\kres(y_1)=K=$ cuerpo de fracciones de~$V$. Sea~$K_s$ la clausura separable
de~$K$, sea~$\overline K$ su clausura algebraica y, para todo subanillo~$W$ de
$\overline K$ que contenga a~$V$, pongamos~$X_W=X\otimes_VW$. En particular,
se tiene
$$
X_V=\ifthenelse{\boolean{orig}}{V}{X}, \qquad X_K=X_1, \qquad X_{\overline K}=\overline X_1.
$$
Por otra parte, el morfismo canónico~$\overline X_1=X_{\overline K}\to X_{K_s}$
induce un isomorfismo en los grupos fundamentales (\Ref{IX.4.11}); por
consiguiente, teniendo en cuenta el isomorfismo~\Ref{X.2.1}
$\pi_1(X_0)\to\pi_1(X)$, quedamos reducidos a estudiar el homomorfismo
sobreyectivo
$$
\pi_1(X_{K_s})\to\pi_1(X)
$$
asociado
\marginpar{281}
al morfismo canónico~$X_{K_s}\to X$. Determinar el núcleo de este último
equivale a resolver el problema siguiente: \emph{se tiene un recubrimiento
principal conexo~$Z_{K_s}$ de~$X_{K_s}$, de grupo~$G$} (asociado, pues, a un
homomorfismo de~$\pi_1(X_{K_s})$ en~$G$); \emph{determinar en qué condiciones
es isomorfo a la imagen inversa de un recubrimiento principal~$Z$ de~$X$ de
grupo~$G$.}

Observemos primero que~$K_s$ es la unión filtrante creciente de sus
subextensiones finitas~$K'$ sobre~$K$ y que, por tanto,~$Z_{K_s}$ es isomorfo a
la imagen inversa de un recubrimiento principal~$Z_{K'}$ de~$X_{K'}$ para un
$K'$ conveniente (adviértase, sin embargo, que para~$K'$ fijado,~$Z_{K'}$ no
queda determinado de manera única). Decir que~$Z_{K_s}$ es isomorfo a la
imagen inversa de un recubrimiento principal~$Z$ de~$X$ significa que existe una
subextensión finita~$K''\supset K'$ de~$K_s$ tal que
$Z_{K''}=Z_{K'}\otimes_{K'}K''$ sea isomorfo a~$Z\otimes_VK''$. Para una
subextensión finita~$K'$ de~$K_s$, designemos ahora por~$V'$ el normalizado
de~$V$ en~$K'$, que es un anillo de valoración discreta completo, de cuerpo
residual~$k$. Así, el morfismo canónico~$X_{V'}\to X_V$ induce un isomorfismo
en las fibras sobre los puntos cerrados de~$Y=\Spec(V)$ y
$Y'=\Spec(V')$. Resulta entonces de~\Ref{X.2.1}, aplicado a~$X_V$ y~$X_{V'}$,
que el homomorfismo inducido en los grupos fundamentales
$\pi_1(X_{V'})\to\pi_1(X_V)$ es un isomorfismo; o, de manera equivalente, que
todo recubrimiento principal de~$X_{V'}$ es la imagen inversa de un
recubrimiento principal de~$X_V$, determinado salvo isomorfismo. Esto implica
el lema siguiente:

\begin{lemme}
\label{X.3.7}
Sea~$Z_{K'}$ un recubrimiento principal conexo de~$X_{K'}$ de grupo~$G$, y sea
\ifthenelse{\boolean{orig}}
{$Z_{K}$}
{$Z_{K_s}$}
su imagen inversa sobre~$X_{K_s}$\quoi. Para que este último sea isomorfo a la
imagen inversa de un recubrimiento principal~$Z$ de~$X$, es necesario y
suficiente que exista una extensión finita~$K''\supset K'$ de~$K$ en~$K_s$ tal
que el recubrimiento principal~$Z_{K''}$ de~$X_{K''}$ esté inducido por un
recubrimiento principal de~$X_{V''}$.
\end{lemme}

Supongamos, en particular, que los~$X_{V''}$ sean normales (para ello basta,
por ejemplo, que~$X_0$ sea normal y, a fortiori, que~$X_0$ sea simple,
\cf \Ref{I.9.1}).
\marginpar{282}
Como son conexos, son, pues, irreducibles. Sean~$L$ el cuerpo de funciones
racionales de~$X$ y~$X_{K}$,~$L'$ el de~$X_{V'}$ y~$X_{K'}$, y~$L''$ el de
$X_{V''}$ y~$X_{K''}$\quoi. Bajo las condiciones de~\Ref{X.3.7},~$Z_{K'}$
define una extensión finita separable~$R'$ de~$L'$, y~$Z_{K''}$ define la
extensión~$R''=R'\otimes_{L'}L'' = R'\otimes_{K'}K''$. La condición
considerada en~\Ref{X.3.7} significa también que existe una extensión finita
separable~$K''$ de~$K'$ tal que~$R''=R'\otimes_{K'}K''$ no esté
\emph{ramificada} sobre el esquema normal~$X_{V''}$, de cuerpo
$L''=L'\otimes_{K'}K''$, y no solamente sobre la parte abierta~$X_{K''}$
de~$X_{V''}$.

Supondremos en adelante que~$f\colon X\to Y$ es un morfismo \emph{liso}; por
tanto, los morfismos~$X_{V'}\to\Spec(V')$ son lisos y los esquemas~$X_{V'}$
son \emph{regulares}. Obsérvese que la fibra del punto cerrado de~$\Spec(V')$
en~$X_{V'}$ es irreducible y de codimensión~$1$. Sea~$\goth{o'}$ su anillo
local, que es, pues, un anillo de valoración discreta de cuerpo~$L'$ y de
cuerpo residual isomorfo al cuerpo de funciones racionales de~$X_0$, y por
tanto de la misma característica que~$k$. Definamos análogamente~$\goth{o''}$
en~$L''$; es evidentemente el normalizado de
\ifthenelse{\boolean{orig}}
{$\goth{o''}$}
{$\goth{o'}$}
en~$L''$. Se deduce entonces del teorema de pureza~\Ref{X.3.1} o~\Ref{X.3.3}
que, para que~$R''$ no esté ramificada sobre~$X_{V''}$, es necesario y
suficiente que~$R''$ no esté ramificada sobre~$\goth{o''}$, normalizado de
$\goth{o'}$ en~$L''$.

Observemos ahora que si~$u'$ es un uniformizante de~$V'$, también lo es
de~$\goth{o'}$. Si~$n$ es un entero primo a la característica~$p$ de~$k$,
\ifthenelse{\boolean{orig}}
{}
{tras ampliar primero~$K'$ de modo que contenga las raíces $n$-esimas de la
unidad,}
y si se toma
\ifthenelse{\boolean{orig}}
{$K''=K'[t]/(t^n-u)$}
{$K''=K'[t]/(t^n-u')$},
entonces
\ifthenelse{\boolean{orig}}
{$K'$}
{$K''$}
es una extensión galoisiana finita de~$K'$ y~$L''$ es isomorfo a
$L'[t]/(t^n-u')$; por tanto, está moderadamente ramificada sobre~$\goth{o'}$ y
su grupo de inercia tiene orden~$n$. Supongamos entonces que~$G$ \emph{tiene
orden primo a~$p$}; esto implica que~$R'$ está moderadamente ramificada
sobre~$\goth{o'}$. Tomemos para~$n$ un múltiplo, primo a~$p$, del orden del
grupo de inercia de~$R'$ sobre~$\goth{o'}$ (por ejemplo,
$n=[G\colon e]$). Aplicando el
\ifthenelse{\boolean{orig}}
{lema de \textsc{Abhyankar}}
{lema de \textsc{Abhyankar}}~\Ref{X.3.6},
se ve que se satisface la condición considerada en~\Ref{X.3.7}.

Esto demuestra el teorema siguiente:

\begin{theoreme}
\label{X.3.8}
Sean~$f\colon X\to Y$ un morfismo propio y liso, de fibras
\marginpar{283}%
geométricamente conexas, con~$Y$ localmente noetheriano, e~$y_0$ e~$y_1$ dos
puntos de~$Y$ tales que
\ifthenelse{\boolean{orig}}
{$y_0\in\overline{y_1}$}
{$y_0\in\overline{\{y_1\}}$};
sean~$\overline{X}_0$ y~$\overline{X}_1$ las fibras geométricas
correspondientes y considérese el homomorfismo de especialización
\eqref{X.2.4}~$\pi_1(\overline X_1)\to\pi_1(\overline X_0)$. Este
homomorfismo es sobreyectivo, y todo homomorfismo continuo de
$\pi_1(\overline X_1)$ en un grupo finito~$G$ de orden primo a la
característica~$p$ de~$\kres(y_0)$ proviene de un homomorfismo de
$\pi_1(\overline{X}_0)$ en~$G$.
\end{theoreme}

En otros términos:

\begin{corollaire}
\label{X.3.9}
Si~$\kres(y_0)$ tiene característica cero, el homomorfismo de especialización
es un isomorfismo. Si
\ifthenelse{\boolean{orig}}
{$p>0$}
{$\kres(y_0)$ tiene característica $p>0$},
entonces el núcleo del homomorfismo de especialización está contenido en la
intersección de los núcleos de los homomorfismos continuos de
$\pi_1(\overline X_1)$ en grupos finitos de orden primo a~$p$ (o, de manera
equivalente, en el subgrupo normal cerrado engendrado por un $p$-subgrupo de
Sylow del grupo de tipo galoisiano~$\pi_1(\overline X_1)$). Por consiguiente,
si
\ifthenelse{\boolean{orig}}
{$\pi_1(X_1)^{(p)}$}
{$\pi_1(\overline X_1)^{(p)}$}
designa el grupo cociente de~$\pi_1(\overline X_1)$ por el subgrupo cerrado
precedente, y si se define análogamente~$\pi_1(\overline X_0)^{(p)}$, el
homomorfismo de especialización induce un \emph{isomorfismo}
$$
\pi_1(\overline X_1)^{(p)}\isomto \pi_1(\overline X_0)^{(p)}
$$
\end{corollaire}

Obsérvese que la demostración de~\Ref{X.3.8} es puramente algebraica.
Procediendo como en~\Ref{X.2.6}, se concluye \emph{por vía trascendente}:

\begin{corollaire}
\label{X.3.10}
Sea~$X_0$ una curva propia, lisa y conexa de género~$g$ sobre un cuerpo
algebraicamente cerrado de característica~$p$. Con la notación introducida
en~\Ref{X.3.9}, el grupo~$\pi_1(X_0)^{(p)}$ es isomorfo a~$\Gamma^{(p)}$,
donde~$\Gamma$ es el grupo de tipo galoisiano engendrado por los generadores
$s_i,~t_i$ ($1\le i\le g$), ligados por la relación
$$
(s_1t_1s_1^{-1}t_1^{-1})\cdots (s_gt_gs_g^{-1}t_g^{-1})=1
$$
\end{corollaire}

\begin{remarques}
\label{X.3.11}
\marginpar{284}
Cuando~$\kres(y_0)$ tiene característica cero, el resultado~\Ref{X.3.9} es bien
conocido por vía trascendente. Obsérvese que la demostración de~\Ref{X.3.10}
recurre al teorema de pureza en el caso de característica mixta, pero
solamente para anillos de dimensión~2, caso en el que la demostración de dicho
teorema es fácil y se ha recordado en el texto.
\end{remarques}

\begin{sgabibliography}{0}{X.4}

\bibitem[1]{X.1} K\ptbl \textsc{Kodaira}, On compact analytic surfaces,
\ifthenelse{\boolean{orig}}
{Princeton University Press, 1960.}
{Ann.\ of Math.\ \textbf{71} (1960), p\ptbl 111--152.}

\bibitem[2]{X.2} S\ptbl \textsc{Lang} y
\ifthenelse{\boolean{orig}}{J.P\ptbl}{J-P\ptbl}%
\textsc{Serre}, Sur les rev\^etements non ramifi\'es des vari\'et\'es
alg\'ebriques,
\ifthenelse{\boolean{orig}}
{Amer.\ Journ. of Math., pp\ptbl 319--330, 1957.}
{Amer.\ J.\ Math.\ \textbf{79} (1957), p\ptbl319--330.}

\bibitem[3]{X.3} M\ptbl \textsc{Nagata},
On the purity of branch loci in regular local rings,
\ifthenelse{\boolean{orig}}
{Illinois Journal of Math., p\ptbl 328--333, 1959.}
{Illinois J.\ Math. \textbf{3} (1959), p\ptbl 328--333.}

\bibitem[4]{X.4}
\ifthenelse{\boolean{orig}}{J.P\ptbl}{J-P\ptbl}%
\textsc{Serre}, Rev\^etements ramifi\'es du plan projectif (d'apr\`es
S\ptbl Abhyankar), S\'eminaire Bourbaki, Mai 1960.

\bibitem[5]{X.5} O\ptbl \textsc{Zariski},
On the purity of the branch locus of algebraic functions,
\ifthenelse{\boolean{orig}}
{Proc.\ Nat.\ Acad.\ Sc.\ USA, t\ptbl 44, p\ptbl 791--796, 1958.}
{Proc.\ Nat.\ Acad.\ Sc.\ USA  \textbf{44} (1958), p\ptbl 791--796.}

\bibitem[6]{X.6} O\ptbl \textsc{Zariski},
Algebraic surfaces, Ergebnisse... 1948, Chelsea (New York).
\end{sgabibliography}
